% TikZ 图：电场加速后进入磁场
% Day 29：模型复盘
\begin{tikzpicture}[scale=0.84,>=Latex,line cap=round,line join=round]
  \path[use as bounding box] (-0.35,-0.95) rectangle (8.10,1.75);
  \draw[thick] (0.30,-0.45) -- (0.30,1.05);
  \draw[thick] (2.05,-0.45) -- (2.05,1.05);
  \node[above] at (0.30,1.05) {$+$};
  \node[above] at (2.05,1.05) {$-$};
  \node[above] at (1.18,1.05) {$U$};
  \foreach \y in {-0.18,0.25,0.68}{
    \draw[->,blue,thick] (0.58,\y) -- (1.72,\y);
  }
  % Place the initially slow positive charge inside the accelerating gap.
  % Putting it outside the positive plate would leave it outside the
  % idealized uniform field shown here.
  \fill[red] (0.58,0.25) circle (1.8pt);
  \node[red,below] at (0.58,0.20) {$+q$};
  \draw[->,red,thick,dashed] (0.72,0.25) -- (2.65,0.25);
  \node[red,font=\scriptsize,above] at (2.25,0.25) {$v_0$};
  \draw[thick] (2.55,-0.55) rectangle (7.55,1.30);
  \foreach \x in {3.0,3.7,4.4,5.1,5.8,6.5,7.2}{
    \foreach \y in {-0.25,0.58,1.15}{
      \node[blue] at (\x,\y) {$\times$};
    }
  }
  % For +q moving right in a field into the page, the center lies directly
  % above the entry point.  The red trajectory and radius share this center.
  \coordinate (P) at (2.65,0.25);
  \coordinate (C) at (2.65,0.95);
  \draw[red,very thick,->] (P) arc[start angle=-90,end angle=0,radius=0.70];
  \fill (C) circle (1.3pt);
  \node[font=\scriptsize,anchor=east] at (2.54,0.98) {$C$};
  \draw[orange!85!black,thick] (P) -- (C);
  \node[orange!85!black,font=\scriptsize,anchor=east] at (2.54,0.60) {$r$};
  \node[font=\small] at (1.18,1.48) {电场加速};
  \node[font=\small] at (5.25,1.48) {磁场偏转};
  \node[blue,font=\scriptsize] at (6.40,-0.30) {$\vec B$ 垂直纸面向里};
\end{tikzpicture}
